\documentclass[11pt,a4paper]{article}

\usepackage[margin=2.4cm]{geometry}
\usepackage{graphicx}
\usepackage{booktabs}
\usepackage{xcolor}
\usepackage{caption}
\usepackage{enumitem}
\usepackage{amsmath}
\usepackage{amssymb}
\usepackage[T1]{fontenc}
\usepackage{lmodern}
\usepackage[colorlinks=true,linkcolor=black,urlcolor=blue,citecolor=black]{hyperref}

\captionsetup{font=small,labelfont=bf}
\setlist{itemsep=2pt,parsep=0pt,topsep=3pt}

\emergencystretch=3em
\newcommand{\code}[1]{\texttt{\small\hyphenchar\font=`\-\relax #1}}
% TODO placeholder: every number below that is not yet measured is wrapped in
% \tbd so that an unfilled table is impossible to mistake for a result.
\newcommand{\tbd}{\textcolor{red}{\textbf{TBD}}}

\title{\textbf{Sketch-Prompted VLA}\\[2pt]
\large A Learned Text-Only Baseline:\\
$\pi_{0.5}$-LIBERO on the Three Validation Suites}
\author{Aaron}
\date{TODO: confirm today's date before compiling}

\begin{document}
\maketitle

\begin{abstract}
\noindent
The 114 validation scenes of this project are ambiguous by construction: the
instruction alone does not identify which of several identical candidates is
meant. Until now the only policies run on them have been scripted --- a
sketch-reading oracle and a firewalled text-only guesser --- which establish that
the harness scores correctly but say nothing about how a real vision-language-action
model behaves under that ambiguity. This report supplies that missing measurement.
I evaluate $\pi_{0.5}$-LIBERO, Physical Intelligence's flow-matching VLA
fine-tuned on the four standard LIBERO suites, as a \emph{learned, sketch-free}
baseline: the stock checkpoint, fed the same instruction string every sketch
condition receives, with no modification to the model. I first reproduce the
published benchmark numbers to certify the observation pipeline, then report
success on my own scenes, separating the two causes of the gap --- distribution
shift away from the standard suites, and referential ambiguity within my tiers.
\tbd{} % headline result sentence, written once the numbers exist
\end{abstract}

\section{Why $\pi_{0.5}$, and why nothing is ablated}

$\pi_{0.5}$ consumes a base image, a wrist image, an eight-dimensional
proprioceptive state and a language string, and emits an action chunk. It has no
sketch channel. ``$\pi_{0.5}$ without the sketch'' therefore requires no edit to
the model, the config or the checkpoint: the sketch-free baseline \emph{is} the
stock checkpoint driven by \code{meta['instruction']}. My claim is that this
makes it the right baseline rather than merely a convenient one --- there is no
ablation whose faithfulness could be disputed, and the number it produces is the
number a competent practitioner would get from the strongest openly available
LIBERO policy with no knowledge of this project.

The checkpoint is \code{gs://openpi-assets/checkpoints/pi05\_libero}, served by
openpi's policy server at config \code{pi05\_libero}. Model and simulator are
kept in separate processes communicating over a websocket, which is how openpi
ships it; that split also means the LIBERO environment used here is
byte-identical to the one the scripted policies run in.

\section{Certifying the observation pipeline}

Four preprocessing details --- render resolution, a $180^{\circ}$ rotation, the
wrist camera, and chunked action execution --- each degrade the success rate
silently rather than raising an error. A low number on my scenes is only
interpretable once they are known to be right, so the published suites are
reproduced first.

\begin{table}[h]
\centering
\caption{Reproduction of the published $\pi_{0.5}$-LIBERO numbers. Published
figures are 500 trials per suite; my reproduction uses a reduced trial count,
sufficient to separate a working pipeline from a broken one.}
\begin{tabular}{lccc}
\toprule
Suite & Published (\%) & Observed (\%) & Trials/task \\
\midrule
LIBERO-Spatial & 98.8 & \tbd & \tbd \\
LIBERO-Object  & 98.2 & \tbd & \tbd \\
LIBERO-Goal    & 98.0 & \tbd & \tbd \\
LIBERO-10      & 92.4 & \tbd & \tbd \\
\midrule
Average        & 96.85 & \tbd & \\
\bottomrule
\end{tabular}
\end{table}

\noindent
\tbd{} % one paragraph: did it reproduce, and what was adjusted if not

\subsection{Divergences from openpi's reference loop}

Two, both deliberate and both recorded in \code{run\_config.json}:

\begin{itemize}
  \item \textbf{No settling wait.} openpi steps ten dummy actions at episode
    start to let dropped objects come to rest. My \code{init\_state.npz} was
    pinned \emph{after} settling and is the state the sketch was drawn on, so
    every condition starts from the identical already-settled configuration and
    the wait is set to zero.
  \item \textbf{Render resolution is decoupled from the sketch canvas.} Scenes
    are rendered at 128 for sketching and at 256 for this policy, from the same
    BDDL and the same restored simulator state. The restored state carries no
    camera information, so the initial physical configuration is identical
    across the two --- which is what makes the conditions comparable.
\end{itemize}

\section{Results on the validation suites}

\begin{table}[h]
\centering
\caption{$\pi_{0.5}$-LIBERO, text-only, on the 114 validation scenes. Sustained
success is the harness's own criterion: the BDDL goal predicate holding over a
five-step window.}
\begin{tabular}{lccccc}
\toprule
Suite & Scenes & Rollouts & Sustained success (\%) & Correct object (\%) & Correct destination (\%) \\
\midrule
Spatial & 38 & \tbd & \tbd & \tbd & \tbd \\
Object  & 38 & \tbd & \tbd & \tbd & \tbd \\
Goal    & 38 & \tbd & \tbd & \tbd & \tbd \\
\midrule
All     & 114 & \tbd & \tbd & \tbd & \tbd \\
\bottomrule
\end{tabular}
\end{table}

\subsection{Separating distribution shift from ambiguity}

The gap between the published average and the figure above has two causes, and
conflating them would overstate what a sketch can recover.

\begin{table}[h]
\centering
\caption{Success by tier. The control tier carries an unambiguous instruction
with a single candidate; it is the honest ceiling for a text-only policy on
these scenes, and the distance from it to the published average is distribution
shift. The distance from control to the other tiers is referential ambiguity ---
the part a sketch prompt is meant to address.}
\begin{tabular}{lccc}
\toprule
Tier & Scenes & Sustained success (\%) & Correct object (\%) \\
\midrule
Control       & \tbd & \tbd & \tbd \\
Other tiers   & \tbd & \tbd & \tbd \\
\bottomrule
\end{tabular}
\end{table}

\noindent
\tbd{} % the interpretation: how much of the gap is which

\subsection{Failure modes}

The harness distinguishes ``grasped the wrong object'' from ``failed to
manipulate at all'' via \code{correct\_instance\_grasped} and
\code{grasped\_any}. That distinction is what makes the headline number mean
something: a policy that cannot grasp in these scenes at all is evidence about
distribution shift, whereas a policy that grasps confidently but chooses the
wrong candidate is evidence about ambiguity, and only the second is the problem
this project claims to solve.

\begin{figure}[h]
\centering
% \includegraphics[width=0.9\textwidth]{figures/fig_tier_breakdown.png}
\caption{TODO: success and failure-mode breakdown by tier.}
\end{figure}

\section{Limitations}

\begin{itemize}
  \item $\pi_{0.5}$-LIBERO was fine-tuned on the four standard suites; my scenes
    are synthesised BDDL with novel arrangements. Its performance here is not a
    measure of the model's quality, and no attempt is made to tune it.
  \item The instruction strings in the non-control tiers are ambiguous by
    design. A perfect policy cannot exceed chance on the referent it is not told
    about, so a portion of the failure rate is irreducible for \emph{any}
    text-only policy --- that portion is precisely what the tier comparison is
    meant to expose.
  \item \tbd{} % anything found during the run
\end{itemize}

\section{Reproducing this}

\begin{itemize}
  \item Brief: \code{prompt\_pi05\_baseline.md}.
  \item Policy: \code{scripts/pi05\_policy.py}; harness path
    \code{scripts/rollout\_sketch\_wsl.py --policy pi05}.
  \item Run artefacts: \code{outputs/rollouts/pi05\_baseline/}
    (\code{results.csv}, \code{run\_config.json}, \code{openpi\_reference.json}).
  \item openpi at \code{main} (not \code{kevin/pi05-support}, which predates the
    merge); git SHA recorded in \code{openpi\_reference.json}.
\end{itemize}

\end{document}
